\documentclass[noauthor]{ximera}

\title{Week 7}
\author{Math 1193 Design Team}

\begin{document}
\begin{abstract}
    Solutions for Week 7.
\end{abstract}
\maketitle

\begin{problem}
Find the following without a calculator. Write your answer as a 
mixed number (for example, \(6 \frac{3}{11}\)). Compare your method 
with your group.

\begin{enumerate}
    \item \(16 \div 3\)
    \item \(242 \div 9\)
    \item \(242 \div 19\)
\end{enumerate}

\begin{explanation}
It's a good idea to remind yourself how long division works.
\begin{enumerate}
    \item \(5 \frac{1}{3}\)
    \item \(26 \frac{8}{9}\)
    \item \(12 \frac{14}{19}\)
\end{enumerate}
\end{explanation}
\end{problem}

\begin{problem}
In the equation \(34 \div 5 = 6 + \frac{4}{5}\), which number
is the divisor, which is the dividend, which is the quotient, and 
which is the remainder?

\begin{explanation}
    Divisor - 5

    Dividend - 34 

    Quotient - 6

    Remainder - 4
\end{explanation}
\end{problem}

\begin{problem}
Check that \(\frac{x^2 + x + 4}{x - 2} = x + 3 + \frac{10}{x - 2}\) by
rewriting the right-hand side using a common denominator. 

Identify the divisor, dividend, quotient, and remainder.

\begin{explanation}
We can rewrite the right-hand side as 
\begin{align*}
    x + 3 + \frac{10}{x - 2} & = \left(\frac{x - 2}{x - 2}\right) \frac{x + 3}{1} + \frac{10}{x - 2} \\
    & = \frac{(x - 2)(x + 3)}{x - 2} + \frac{10}{x - 2} \\
    & = \frac{(x - 2)(x + 3) + 10}{x - 2} \\
    & = \frac{x^2 - 2x + 3x - 6 + 10}{x - 2} \\
    & = \frac{x^2 + x + 4}{x - 2},
\end{align*}
so the two sides are equal. 

Divisor - \((x - 2)\)

Dividend - \((x^2 + x + 4)\)

Quotient - \((x + 3)\)

Remainder - \(10\)
\end{explanation}
\end{problem}

\begin{problem}
\begin{enumerate}
    \item Check that \(x^2 - 1 = (x + 2)(x - 2) + 3\) by rewriting the right-hand
    side. You'll need to multiply and collect like terms.
    \item What is the result when plugging in \(x = 2\) to the right-hand side?
    \item Therefore, what is the result when plugging in \(x = 2\) to the left-hand
    side? You shouldn't have to do any extra work here, since the
    left- and right-hand sides are equal.
\end{enumerate}

\begin{explanation}
    \begin{enumerate}
        \item We can rewrite the right-hand side as
        \begin{align*}
            (x + 2)(x - 2) + 3 & = (x^2 + 2x - 2x - 4) + 3 \\
            & = x^2 - 1
        \end{align*}
        which is what we expected.

        \item If we plug in \(x = 2\) to the right-hand side, we obtain \((2 + 2)(2 - 2) + 3 = 4 \cdot 0 + 3 = 3\).
        \item If we plug in \(x = 2\) to the left-hand side, we also obtain 3, since it equals the right-hand side.
    \end{enumerate}
\end{explanation}
\end{problem}

\begin{problem}
Rational functions are fractions where the numerator and
denominator are polynomials.

Is \(f(x) = \frac{2}{2 - x} - \frac{1}{x - 4}\) a rational function?
It's a \textbf{sum} of two fractions of polynomials, but can it be
written as just \textbf{one} fraction of polynomials?

Write \(f(x)\) as one fraction of polynomials.

\begin{explanation}
We can rewrite the function as 
\begin{align*}
    \frac{2}{2 - x} - \frac{1}{x - 4} & = \left(\frac{x - 4}{x - 4}\right) \frac{2}{2 - x} - \frac{1}{x - 4} \left(\frac{2 - x}{2 - x}\right) \\
    & = \frac{(x - 4)(2)}{(x - 4)(2 - x)} - \frac{(2 - x)}{(x - 4)(2 - x)} \\
    & = \frac{(x - 4)(2) - (2 - x)}{(x - 4)(2 - x)}  \\
    & = \frac{2x - 8 - 2 + x}{(x - 4)(2 - x)} \\
    & = \frac{3x - 10}{(x - 4)(2 - x)},
\end{align*}
so yes, this can be written as one fraction of polynomials.
\end{explanation}
\end{problem}

\begin{problem}
Let \(C\) be the function that gives the average cost per book 
\(C\left(x\right)\), in dollars, when using an online store to print
\(x\) copies of a self-published paperback book. Here is a graph of 
\(C\left(x\right)=\frac{120+4x}{x}\). 

\begin{center}
    \includegraphics[width=10cm]{bookselling.png}
\end{center}

\begin{enumerate}
    \item What is the approximate cost per book when 50 books are 
    printed? 100 books?
    \item The author plans to charge \$8 per book. About how many
    should be printed to make a profit?
    \item What is the value of \(C(x)\) when \(x = \frac{1}{2}\)? how
    does this relate to the context?
    \item What does the end behavior of the function \(C\) say about
    the context. Write down your answer, then share with a partner.
\end{enumerate}

\begin{explanation}
\begin{enumerate}
    \item To find the cost per book, we plug the number of books printed into  our function \(C\).
    
    \(C(50) = \frac{320}{50} = 6.4\), so \$6.40 per book, and \(C(100) = \frac{520}{100} = 5.2\), so \$5.20 per book.

    \item If the author wants to make a profit, they need to make back at least \$8 per book, so the cost per book has to be no more than
    \$8. The cost is \$8 per book when \(C(x) = 8\), which is an equation we can solve:
    \begin{align*}
        \frac{120 + 4x}{x} & = 8 \\
        120 + 4x & = 8x \\
        120 & = 4x \\
        30 & = x
    \end{align*}
    Thus, the author makes a profit when they make more than 30 books.

    \item \(C\left(\frac{1}{2}\right) = \frac{120 + 4\left(\frac{1}{2}\right)}{\frac{1}{2}} = \frac{120 + 2}{1} \cdot \frac{2}{1} = 122 \cdot 2 = 244\)
    
    This means that when half a book is produced, the cost per book is \$244. That doesn't make much sense, since people usually don't produce half books, but 
    it could be the cost per book when the book is half as long.

    \item The function's graph has a horizontal asymptote at  \(y = 4\), meaning as \(x\) approaches positive infinity, \(C(x)\) approaches 4.
    The more books you produce, the closer the cost per book gets to \$4. 
\end{enumerate}
\end{explanation}
\end{problem}

\begin{problem}
For each of the following functions, solve the equation \(f(x) = 0\)
and evaluate \(f(0)\). What do these answers tell you about the graph 
of \(f\)?

\begin{enumerate}
    \item \(f(x) = \frac{x^2 - 1}{x + 2}\)
    \item \(f(x) = \frac{x + 1}{x^2 - 1}\)
    \item \(f(x) = \frac{1}{x - 3} + \frac{2}{x + 1}\)
    \item \(f(x) = \frac{2x + 1}{x}\)
\end{enumerate}

\begin{explanation}
Remember that \textbf{fractions are only 0 when their numerators are 0}, so to solve \(f(x) = 0\), we can just set the numerator equal to 0.
\begin{enumerate}
    \item 
    \begin{align*}
        x^2 - 1 & = 0 \\
        (x - 1)(x + 1) & = 0
    \end{align*}
    so by setting each factor equal to 0, we obtain \(x - 1 = 0\) and \(x + 1 = 0\), so \(x = -1, 1\) are the solutions to the equation \(f(x) = 0\).

    \(f(0) = \frac{0^2 - 1}{0 + 2} = \frac{-1}{2}\)

    \item
    \begin{align*}
        x + 1 & = 0 \\
        x & = -1
    \end{align*}
    so \(x = -1\) is the solution to the equation \(f(x) = 0\).

    \(f(0) = \frac{0 + 1}{0^2 - 1} = \frac{1}{-1} = -1\)

    \item It's helpful to rewrite this as one fraction:
    \begin{align*}
        \frac{1}{x - 3} + \frac{2}{x + 1} & = \left(\frac{x + 1}{x + 1}\right)\frac{1}{x - 3} + \frac{2}{x + 1} \left(\frac{x - 3}{x - 3}\right) \\
        & = \frac{x + 1}{(x + 1)(x - 3)} + \frac{2(x - 3)}{(x + 1)(x - 3)} \\
        & = \frac{(x + 1) + 2(x - 3)}{(x + 1)(x - 3)} \\
        & = \frac{x + 1 + 2x - 6}{(x + 1)(x - 3)} \\
        & = \frac{3x - 5}{(x + 1)(x - 3)} \\
    \end{align*}

    So \(f(x) = 0\) can be solved by solving \(3x - 5 = 0\), meaning \(x = \frac{5}{3}\).

    Finally, \(f(0) = \frac{3 \cdot 0 - 5}{(0 + 1)(0 - 3)} = \frac{-5}{-3} = \frac{5}{3}\).

    \item 
    \begin{align*}
        2x + 1 & = 0 \\
        2x & = -1 \\
        x & = -\frac{1}{2}
    \end{align*}
    This is the solution to \(f(x) = 0\). 

    \(f(0)\) is undefined, since there cannot be a 0 in the denominator.
\end{enumerate}
\end{explanation}
\end{problem}

\end{document}
